%% Header mit Deklarationen
\documentclass[%
paper=a4,      % alle weiteren Papierformat einstellbar
fontsize=11pt, % Schriftgr��e (12pt, 11pt (Standard))
BCOR1cm,       % Bindekorrektur, bspw. 1 cm
DIV15,         % f�hrt die Satzspiegelberechnung neu aus s. scrguide 2.4
twoside,       % Doppelseiten
headsepline,   %
headings=openright, % Kapitel nur rechts beginnen
%biblography=totoc, % Literaturverzeichnis einf�gen bibtotocnumbered: nummeriert
parskip=half,  % Europ�ischer Satz mit Abstand zwischen Abs�tzen
chapterprefix, % Kapitel anschreiben als Kapitel
headsepline,   % Linie nach Kopfzeile
titlepage,     %
numbers=noenddot,
%draft	       % zeigt �berlange Zeilen an
]{scrreprt}

\usepackage{pdfpages}       % Titelseite hat ein anderes Layout. Sie wird 
                            % separat erzeugt und hier eingef�gt
\usepackage[T1]{fontenc}
\usepackage[latin1]{inputenc}  % Zeichencodierung
\usepackage[ngerman, english]{babel} % Worttrennung nach neuer Rechtschreibung
\usepackage{ellipsis}       % Leerraum um Auslassungspunkte
\usepackage{fixltx2e}       % Fehlerkorrektur Zeichens�tze
\usepackage{xspace}         % f�ge evtl. notwendiges Leerzeichen hinzu (\xspace)

%\usepackage{mathptmx}           % Times + passende Mathefonts
\usepackage{mathpazo}           % Palatino + passende Mathefonts
\usepackage[scaled=.92]{helvet} % skalierte Helvetica als \sfdefault
\usepackage{courier}            % Courier als \ttdefault

\usepackage{graphicx}    % Einbindung von Grafiken
\graphicspath{{bilder/}} % Unterverzeichnis, in dem Grafiken abgelegt werden
\usepackage{listings}    % Listenausgabe externer Dateien

\usepackage{float}      % Paket zum Erweitern der Floatumgebungen
\usepackage[figuresright]{rotating}   % Rotieren von Objekten
%\usepackage{hvfloat}
\usepackage{array}      % Paket zum Erweitern der Tabelleneigenschaften
\usepackage{booktabs}   % Paket f�r sch�nere Tabellen

\usepackage{amsmath}    % erweiterte Mathematik-Umgebungen
\usepackage{amssymb}

% Einstellungen f�r das Literaturverzeichnis
\usepackage[round]{natbib}
\setlength{\bibsep}{0.5\baselineskip}
\setlength{\bibhang}{1cm}
\bibliographystyle{agsm}

% Andere Schriftarten in Koma-Script
\setkomafont{sectioning}{\normalfont\bfseries}
\setkomafont{captionlabel}{\rmfamily\bfseries\small}
\setkomafont{caption}{\mdseries\itshape\small}
\setkomafont{pagehead}{\normalfont\itshape} % Kopfzeilenschrift
\setkomafont{descriptionlabel}{\normalfont\bfseries}

% Kopf und Fu�zeilen
\usepackage[automark]{scrpage2}

% Literaturverzeichnis-Stil
\bibliographystyle{plain}

% weitere Einstellungen
\tolerance=200               % �bervolle Zeile vermeiden
\emergencystretch=3em

\clubpenalty=10000           % 'Schusterjungen' und 'Hurenkinder' vermeiden
\widowpenalty=10000 
\displaywidowpenalty=10000

\parindent 0pt               % Einzug zu Absatzbeginn festlegen

\setcapindent{1em}           % Zeilenumbruch bei Bildbeschreibungen.

\setcounter{secnumdepth}{3}  % Strukturiertiefe bis subsubsection{} m�glich
\setcounter{tocdepth}{3}     % Dargestellte Strukturiertiefe im Inhaltsverzeichnis

% Korrekturversion mit 1.5-fachem Zeilenabstand im Hauptteil:
\newif\ifiscorrect
%\iscorrecttrue   % Korrekturversion
\iscorrectfalse % keine Korrekturversion

%% Eigene Definitionen:

% Einheiten:
\def\ut#1{\ensuremath{\,\mathrm{#1}}}

% Operatoren:
\def\grad{\ensuremath{\mathop{\mathrm{grad}}\nolimits}}
\def\transp#1{\ensuremath{{#1}^\mathsf{T}}}  % transpose
\def\const{\ensuremath{\mathop{\mathrm{const.}}\nolimits}}

% Formelzeichen:
\def\vec#1{\ensuremath{\mathbf{#1}}}
\def\matr#1{\ensuremath{\mathbf{#1}}}

% Hack, um ein zus�tzliches Leerzeichen nach \input zu entfernen:
\def\myinput#1{%
  \endlinechar=-1 % kein Zeilenabschlusszeichen
  \input #1\relax
  \endlinechar `\^^M % Zeilenabschluss = Zeilenvorschub
}


\begin{document}

% R�mische Nummerierung f�r Sonderseiten, wie Verzeichnisse und Anhang
\pagenumbering{Roman}
\pagestyle{plain}

%% Titelblatt
\includepdf[pages=1]{titelseite}
\cleardoublepage

%% Eidesstattliche Erkl�rung
% Die eidesstattliche Erkl�rung mit Unterschrift
\chapter*{Erkl�rung der Urheberschaft}

Ich erkl�re hiermit an Eides statt, dass ich die vorliegende Arbeit
ohne Hilfe Dritter und ohne Benutzung anderer als der angegebenen
Hilfsmittel angefertigt habe; die aus fremden Quellen direkt oder
indirekt �bernommenen Gedanken sind als solche kenntlich gemacht. Die
Arbeit wurde bisher in gleicher oder �hnlicher Form in keiner anderen
Pr�fungsbeh�rde vorgelegt und auch noch nicht ver�ffentlicht.


\vspace{4cm}

\hspace{2cm} Ort, Datum \hfill Unterschrift \hspace{2cm}

\cleardoublepage

%% Abstract
\begin{abstract}
In diesem kurzen Text erfolgt die Zusammenfassung der Arbeit oder -- auf Englisch -- das Abstract.

\end{abstract}
\cleardoublepage

%% Verzeichnisse
% Inhalt Kopf-/Fu�zeile
\pagestyle{scrheadings}
\clearscrheadfoot            % Standardkram wegwerfen
\ohead[\pagemark]{\pagemark} % oben au�en Seitenzahl 
%
% Inhaltsverzeichnis
%
\tableofcontents

%
% Abbildungsverzeichnis
%
\listoffigures

%
% Tabellenverzeichnis
%
\listoftables


% Merke mir die r�mische Seitenzahl in 'roemisch' und setzte Nummeriernung 
% auf arabisch f�r die eigentlichen Kapitel
\cleardoublepage %\newpage
\newcounter{roemisch}
\setcounter{roemisch}{\value{page}}
\pagenumbering{arabic}

\ifiscorrect\linespread{1.5}\selectfont% Zeilenabstand: 1 1/2 (f�r bessere Korrektur)
\else\fi

%% Die einzelnen Kapitel
% Kopfzeile: links Kapitel, rechts Sektion
\clearscrheadfoot            % Standardkram wegwerfen
\ohead[\pagemark]{\pagemark} % oben au�en Seitenzahl 
\ihead{\headmark}            % oben innen automatischer Abschnittsname
%\automark[]{section}
\chapter{Erstes Kapitel}
\section{Erster Abschnitt}

Hier beginnt der erste Satz, der dann gleich mit einem Bild weitergef�hrt wird
\begin{figure}[htbp]
\centering
\includegraphics[width=0.7\textwidth]{logo_gi} % Datei in "bilder/" bei LaTeX: eps, bei PDFLaTeX: jpg (o.�.) 
\caption{Titel der Abbildung} 
\label{fig:bild1}
\end{figure}

In der Abbildung~\ref{fig:bild1}\footnote{vgl. Zitat A \citep{Mustermann:2003}} % referenz A muss im Verzeichnis "literatur/" in der Datei "bib.bib" enthalten sein, f�r weitere Befehle f�rs Einf�gen von Zitaten oder Quellenverweisen siehe die Dokumentation zum Paket "natbib"
ist zu sehen, dass ...

... Einheiten werden mit dem korrekten Abstand zur Zahl gesetzt: 123\ut{m}.

% Neue Seite
\newpage
\section{Und n�chster Abschnitt}
Eine neue Seite, um auchmal die Kopfzeile zu sehen, da sie auf Seiten mit Kapitelanfang nicht erscheinen...

\chapter{Zweites Kapitel}
\section{Neuer Abschnitt}

Hier f�ge ich mal eine Tabelle ein
\begin{table}[htbp]
\centering
\begin{tabular}{lrcl}
\toprule
SpalteA & SpalteB & SpalteC & SpalteD \\
\midrule
InhaltA1 & InhaltB1 & InhaltC1 & InhaltD1 \\
InhaltA2 & InhaltB2 & InhaltC2 & InhaltD2 \\
InhaltA3 & InhaltB3 & InhaltC3 & InhaltD3 \\
\bottomrule
\end{tabular}
\caption{Beispiel einer Tabelle}
\label{tab:tabelle1}
\end{table}

Wie man in der Tabelle~\ref{tab:tabelle1} sehen kann ...

\chapter{Drittes Kapitel}
\section{Absatz}

Hier mal eine Aufz�hlung
\begin{enumerate}
 \item erster Punkt
 \item noch ein Punkt
 \item letzter Punkt
\end{enumerate}

Super, oder?

\chapter{Viertes Kapitel}
\section{Eine Sektion}

Hier mal eine Auflistung von Elementen
\begin{itemize}
 \item erstes Element
 \item zweites Element
 \item noch ein Element
\end{itemize}

Und Schluss mit der Vorlage ...


\ifiscorrect\linespread{1.0}\selectfont% Zeilenabstand wieder auf 1 zur�ck
\else\fi

% Setze Numerierung wieder auf r�misch zur�ck und setzte von oben fort
% Wert ist demnach der von 'roemisch'
\newpage
\pagenumbering{Roman}
\setcounter{page}{\value{roemisch}}

%% Literaturverzeichnis
\bibliography{literatur/bib}

%% Appendix, falls vorhanden
\appendix
\chapter[Anhang]{}

Und das hier ist noch der Anhang...


\end{document}
